\newpage
\section{FUTURE ENHANCEMENTS AND CONCLUSION}

\subsection{Future Enhancements}

Although the final prototype successfully demonstrates the feasibility of RL-based
antenna alignment, several technical improvements can further enhance system accuracy,
robustness, and deployability.

\subsubsection{Mechanical and Hardware Improvements}

The current prototype performs well for proof-of-concept validation, but future hardware
revisions can improve long-term stability and measurement repeatability. In particular,
mechanical backlash in the pan axis and minor vibration after step movement can be reduced
through a tighter coupling mechanism, higher-quality bearings, and a more rigid antenna
mount. Similarly, the custom slip-ring assembly can be refined using improved contact
pressure control and better conductive surface finishing to reduce intermittent contact
noise during long-duration rotation.

A higher-gain directional antenna and a calibrated RSSI front-end may also improve the
signal-to-noise ratio of the measurements. This would be especially useful in indoor and
obstructed conditions where the system must discriminate between several nearby local RSSI
peaks.

\subsubsection{Algorithmic Improvements}

The present work uses a compact Q-learning controller because it is simple, interpretable,
and easy to deploy on a low-cost microcontroller. However, future work can explore more
advanced learning strategies to improve adaptation under dynamic channel conditions.
Possible improvements include:
\begin{itemize}
    \item \textbf{Online policy adaptation:} Allow limited on-device updating of selected Q-values or policy parameters when the propagation environment changes significantly.
    \item \textbf{Function approximation:} Replace the tabular policy with a lightweight neural approximation or compressed value model to support finer angular resolution without a large memory increase.
    \item \textbf{Hybrid search strategy:} Combine a rapid coarse scan with RL-based refinement to further reduce convergence time in weak-signal conditions.
    \item \textbf{Temporal state features:} Include short RSSI history, packet loss rate, or moving-average slope so that the agent can better distinguish between stable improvement and temporary fading.
\end{itemize}

These improvements would allow the controller to respond more effectively to non-stationary
channels, moving users, and changing obstruction patterns.

\subsubsection{Embedded Software Enhancements}

The firmware currently uses a deterministic blocking control loop, which was appropriate
for experimental validation. In future versions, the software can be improved by adopting
interrupt-assisted sampling or an RTOS-based scheduling model. This would enable parallel
handling of RSSI acquisition, telemetry logging, and actuator control, thereby improving
responsiveness without compromising measurement integrity.

Additional software enhancements may include automatic experiment logging to onboard
storage, wireless transfer of recorded trial data, remote parameter tuning, and a simple
user interface for selecting alignment modes and calibration routines.

\subsubsection{System-Level Extensions}

The proposed system can be extended beyond the present laboratory-scale implementation.
A multi-node version could coordinate several receivers or repeaters simultaneously,
allowing distributed beam steering in wireless mesh or relay networks. The alignment logic
could also be integrated with link-layer performance indicators such as throughput, packet
error rate, or signal-to-noise ratio, thereby enabling the controller to optimize overall
communication quality rather than RSSI alone.

\begin{table}[H]
\centering
\caption{Recommended Future Technical Enhancements}
\label{tab:future_enhancements}
\renewcommand{\arraystretch}{1.4}
\begin{tabular}{|p{3.7cm}|p{4.2cm}|p{4.2cm}|}
\hline
\textbf{Enhancement Area} & \textbf{Proposed Improvement} & \textbf{Expected Benefit} \\
\hline
Mechanical design & Reduce backlash, improve bearing support, refine slip-ring contact & More accurate positioning and lower measurement disturbance \\
\hline
RF sensing & Use higher-gain antenna and calibrated RSSI acquisition & Improved resolution between candidate directions \\
\hline
Learning algorithm & Introduce hybrid RL or lightweight function approximation & Better generalization and finer control resolution \\
\hline
Firmware architecture & Move toward event-driven or RTOS-assisted execution & Faster sensing-actuation coordination and improved scalability \\
\hline
System integration & Add telemetry, remote tuning, and persistent experiment logging & Easier deployment, monitoring, and maintenance \\
\hline
Network-level intelligence & Optimize with throughput or SNR in addition to RSSI & Better end-to-end communication performance \\
\hline
\end{tabular}
\end{table}

\subsection{Conclusion}

This project developed and validated an autonomous directional antenna alignment system
based on reinforcement learning and low-cost embedded hardware. The work combined a custom
mechanical platform, ESP32-S3-based sensing and actuation, offline Q-learning, and both
simulation and practical evaluation into a complete end-to-end prototype.

The final results show that the proposed RL-based controller outperformed the selected
baseline methods in the most important practical metrics. Compared with exhaustive
scanning, the RL controller achieved substantially lower convergence time while also
producing stronger final RSSI in the tested scenarios. Compared with hill-climbing, the
RL method showed superior robustness, better convergence reliability, and greater immunity
to local maxima, particularly in multipath-rich indoor trials and partially obstructed
outdoor trials.

The practical experiments at 10~m and 20~m, in both indoor and outdoor environments,
confirm that the RL-based approach is not only theoretically promising but also feasible
for real embedded implementation. The trained Q-table remained compact enough for direct
microcontroller deployment, and the resulting control loop was computationally lightweight
while still delivering a measurable performance gain.

Overall, the project demonstrates that reinforcement learning can provide an effective and
practical solution for autonomous antenna alignment in low-cost wireless systems. The
prototype and results presented in this report establish a strong foundation for future
work on adaptive beam steering, intelligent wireless links, and embedded RL-driven control
systems.
